\subsection{Агроном-любитель}

\subsubsection{Описание решения}

Переформулировав условие задачи, в массиве нужно найти наиболее длинный отрезок, в котором отсутствуют три одинаковых числа подряд.

Заметим, что для достижения максимальной длины отрезка необходимо, чтобы искомый отрезок начинался как можно раньше и заканчивался как можно позже. Другими словами, если \( i \)-ый отрезок можно расширить без нарушения условия --- мы расширяем.

Рассмотрим ситуацию, когда в массиве в первый раз встретились три и более одинаковых числа подряд: \( x_i=x_{i+1}=x_{i+2}=\dots \). Следуя наблюдению выше, первый отрезок должен сделать \( x_{i+1} \) своим последним элементом. Тогда следующий отрезок должен начинаться с \( x_{i+1} \) элемента. Для следующей встречи рассуждения аналогичны, как если бы исходный массив начинался с \( x_i \) элемента.

Учитывая, что отрезки пересекаются ровно по одному элементу, можно пройтись по всему массиву один раз и вести учёт длины каждого отрезка, отдельно сохраняя наибольший из них.

\subsubsection{Асимптотическая сложность}

Временная сложность решения --- линейная \( \mathcal{O}(N) \). В алгоритме для каждого элемента массива выполняется константное количество операций.

\subsubsection{Листинг кода}

\begin{lstlisting}
#include <iostream>

struct segment {
  int l;
  int r;
};

int main() {
  int N;
  std::cin >> N;
  struct {
    short occur = 0;
    long item = -1;
  } history;
  segment ans = {0, 0};
  segment curr = {0, 0};
  for (int i = 0; i < N; i++) {
    long a;
    std::cin >> a;
    if (history.item != a) {
      history.item = a;
      history.occur = 1;
    } else {
      history.occur++;
    }
    if (history.occur == 3) {
      if (ans.r - ans.l < curr.r - curr.l) {
        ans = curr;
      }
      history.occur--;
      curr.l = i - 1;
    }
    curr.r = i;
  }
  if (ans.r - ans.l < curr.r - curr.l) {
    ans = curr;
  }
  std::cout << ans.l + 1 << " " << ans.r + 1;
  return 0;
}
\end{lstlisting}